\documentclass[10pt,a4paper]{amsart}
\usepackage[margin=19mm]{geometry}
\usepackage{amsmath,amssymb,mathtools}
\usepackage[scr=rsfs,cal=euler]{mathalfa}
\usepackage{xfrac}
\usepackage[unicode,hidelinks,bookmarks=false]{hyperref}
\newcommand{\ie}{i.e.\ }
\newcommand{\cf}{cf.\ }
\newcommand{\resp}{respectively\ }
\newcommand{\Subsection}{\subsection}
\providecommand{\nobreakdash}{\nobreak\hbox{-}}
\setlength{\emergencystretch}{2em}
\multlinegap0pt
\usepackage{bm}
\usepackage{bbm}
\usepackage{dsfont}
\usepackage{xspace}
\usepackage{cancel}
\usepackage{mathtools}
\usepackage{amsthm}
\makeatletter
\@ifundefined{definition}{\newtheorem{definition}{Definition}}{}
\@ifundefined{lemma}{\newtheorem{lemma}{Lemma}}{}
\@ifundefined{proposition}{\newtheorem{proposition}{Proposition}}{}
\@ifundefined{theorem}{\newtheorem{theorem}{Theorem}}{}
\makeatother
\makeatletter
\expandafter\ifx\csname smallbmatrix\endcsname\relax
\newenvironment{smallbmatrix}{\left[\begin{smallmatrix}}{\end{smallmatrix}\right]}
\fi
\makeatother
\title[On the Regularity and Interpolation of Coupled Cluster Ampli]{On the Regularity and Interpolation of Coupled Cluster Amplitudes in Canonical Orbital Basis}
\author{Jonas Beck, Benjamin Stamm}
\date{Source: arXiv:2605.22584v2 (CC BY 4.0)}
\begin{document}
\maketitle

\noindent\emph{Source and licence.} On the Regularity and Interpolation of Coupled Cluster Amplitudes in Canonical Orbital Basis. Version arXiv:2605.22584v2 (CC BY 4.0), distributed under the Creative Commons Attribution 4.0 International licence, as recorded in the arXiv licence metadata for this exact version and shown on the abstract page https://arxiv.org/abs/2605.22584v2. Full rights notice: licenses/arxiv-2605.22584-CC-BY-4.0.txt.\par\smallskip

\noindent\textit{Editorial scope.} Counted unit: the Theorem whose source label is \texttt{analytic\_amplitudes} in the article (source0.tex lines 1007--1017, source label \texttt{analytic\_amplitudes}), delivered together with the article's own complete original proof of it (source0.tex lines 1019--1055, one \texttt{proof} environment of 37 lines). No mathematics is written, repaired or re-derived: statement and proof are the article's own text line for line. Required context retained uncounted: proposition\_excitation\_comm (lines 731--738); trunc\_slater\_basis (lines 758--760); cc\_energy (lines 786--788); cluster\_function (lines 974--976); creation\_annihili\_analytic (lines 978--980); non-degen\_defi (lines 998--1004). Every definition, display and label of those lines is retained unchanged. Omitted from this package and not redistributed: 1--730, 739--757, 761--785, 789--973, 977--977, 981--997, 1005--1006, 1018--1018, 1056--1310. Nothing inside a retained excerpt is omitted or altered: every retained body is delivered exactly as the article's lines are written, so no omission receipt of any kind accompanies this package. Citations: the retained excerpts carry no citation command at all, so no bibliography entry of the article is transcribed and no reference is redistributed. Reference adaptation: none was needed and none was made; every reference inside the delivered unit resolves inside it, so no unresolved reference or citation is produced. Printed slips kept literal: no printed word, symbol, hypothesis or number of the counted statement or of its proof was altered. Layout and class: the article's own class is \texttt{amsart} with packages including macros/preprintStyle, bm, comment, adjustbox, framed, xcolor, amsthm, tikz, mathrsfs, rsfso, dutchcal, mdframed, booktabs, multirow; the wrapper is the lane's pinned \texttt{amsart} editorial wrapper at 10pt on A4, which restates the theorem-like environments the retained text needs and adds the article's own macro definitions that the retained text needs (none), with extra wrapper packages (bm, bbm, dsfont, xspace, cancel, mathtools). The delivered numbering is the unit's own. Assets: no retained span contains a figure environment, a graphics-inclusion command, a PDF-page inclusion, a table or a TikZ picture; no image file is copied into this package, the package declares no asset dependency and no image file is redistributed with it. Licence: version arXiv:2605.22584v2 of this article is distributed under the Creative Commons Attribution 4.0 International licence, as recorded in the arXiv licence metadata for this exact version and shown on the abstract page https://arxiv.org/abs/2605.22584v2. A full rights notice is delivered at licenses/arxiv-2605.22584-CC-BY-4.0.txt. Characters: every retained excerpt is pure ASCII, so the delivered engine, XeLaTeX with its default Latin Modern text font, reports no missing character.
\begin{proposition}\label{proposition_excitation_comm}
    For any $\nu, \eta \in \mathscr{I}$ it holds that 
    \begin{enumerate}
        \item (commutativity) $[\mathscr{X}_\nu,\mathscr{X}_\eta]=\mathscr{X}_\nu \mathscr{X}_\eta- \mathscr{X}_\eta \mathscr{X}_\nu = 0=\mathscr{X}^\dagger_\nu \mathscr{X}^\dagger_\eta- \mathscr{X}^\dagger_\eta \mathscr{X}^\dagger_\nu=[\mathscr{X}^\dagger_\nu,\mathscr{X}^\dagger_\eta]$.
        \item (closedness) either $\mathscr{X}_\nu \mathscr{X}_\eta=0$ or there exists $\upsilon \in \mathscr{I}$ with $\mathscr{X}_\nu \mathscr{X}_\eta=\mathscr{X}_\upsilon$.
        \item (nilpotency)  $\big(\mathscr{X}^\dagger_\nu\big)^2=\mathscr{X}_\nu^2=0$.
    \end{enumerate}
\end{proposition}

\begin{align}\label{trunc_slater_basis}
    \mathcal{S}_\ell(\mathfrak{M}):=\lbrace \Psi_0\rbrace \cup \lbrace \mathscr{X}_\nu \Psi_0\,|\, \nu \in \mathscr{I}_\ell\rbrace.
\end{align}

\begin{align}\label{cc_energy}
     E_{\ell}:=E_{CC}:=\langle \Psi_0, \mathrm{e}^{-T}\mathcal{H}_\mathfrak{M}\mathrm{e}^T\Psi_0\rangle_{\mathrm{L}^2}.
\end{align}

\begin{align}\label{cluster_function}
    T:\Omega\times\mathbb{R}^{|\mathscr{I}_\ell|}\longrightarrow \mathcal{B}(\mathfrak{F}(\mathfrak{h})), \quad (\omega, t)\mapsto \sum_{\nu\in \mathscr{I}_\ell}t_\nu \mathscr{X}_\nu(\omega).
\end{align}

\begin{lemma}\label{creation_annihili_analytic}
    Suppose $\mathfrak{M}(\omega)=\lbrace \psi_i(\omega)\rbrace_{i=1}^N\cup \lbrace \psi_a(\omega)\rbrace_{a=N+1}^{\mathrm{n}_\mathrm{b}}$ is set of real analytic molecular orbitals as Banach space-valued functions. Then the corresponding creation operators $\lbrace a^\dagger_i(\omega)\rbrace_{i=1}^N\cup \lbrace a^\dagger_a(\omega)\rbrace_{a=N+1}^{\mathrm{n}_\mathrm{b}} $ and annihilation operators $\lbrace a_i(\omega)\rbrace_{i=1}^N\cup \lbrace a_a(\omega)\rbrace_{a=N+1}^{\mathrm{n}_\mathrm{b}}$ are real analytic, as well. Moreover, any excitation operator $\mathscr{X}_\nu(\omega)$, for $\nu\in \mathscr{I}$, and the cluster function \eqref{cluster_function} are also real analytic. 
\end{lemma}

\begin{definition}\label{non-degen_defi} Let $\ell\le N$ be a fixed excitation level, $\mathfrak{M}(\omega)= \lbrace \psi_i(\omega)\rbrace_{i=1}^N\cup \lbrace \psi_a(\omega)\rbrace_{a=N+1}^{\mathrm{n}_\mathrm{b}}$ be a set of molecular orbitals and $\mathcal{Q}$ the coupled cluster function with respect to the truncated Slater basis $\mathcal{S}_\ell(\mathfrak{M}(\omega))$ (see \eqref{trunc_slater_basis}). 
    A root $(\omega^*, t^*)$ of the parameter dependent coupled cluster function $\mathcal{Q}$ is a non-degenerate root if for $\mathfrak{V}_\ell:=\mathrm{span}_{\mathbb{C}}\,\mathcal{S}_\ell(\mathfrak{M}(\omega^*))\setminus  \lbrace \Psi_0(\omega^*) \rbrace$ 
    \begin{align}\label{non-degen-assum}
        E_{CC} \notin \sigma\big(P_{\mathfrak{V}_\ell}\mathrm{e}^{-T(\omega^*,t^*)}\mathcal{H}_{\mathfrak{M}}(\omega^*)\mathrm{e}^{T(\omega^*,t^*)}P_{\mathfrak{V}_\ell}\big),
    \end{align}
  where $E_{CC}$ is the coupled cluster energy at point $(\omega^*,t^*)$ (see \eqref{cc_energy}) and $P_{\mathfrak{V}_\ell}$ the orthogonal projector onto $\mathfrak{V}_\ell$.
\end{definition}

\begin{theorem}\label{analytic_amplitudes}

    Let $\Omega\subset \mathbb{R}^{3M}$ be an open subset of the nuclear coordinate space and $\ell\le N$ be a fixed excitation level. Furthermore, let  $\mathfrak{M}(\omega)= \lbrace \psi_i(\omega)\rbrace_{i=1}^N\cup \lbrace \psi_a(\omega)\rbrace_{a=N+1}^{\mathrm{n}_\mathrm{b}}$ be a set of real analytic Banach space valued molecular orbitals and $\mathcal{S}_\ell(\mathfrak{M}(\omega))$ the corresponding truncated set of Slater determinants. If $t^* \in \mathbb{R}^{|\mathscr{I}_\ell|}$ is the cluster amplitude for a certain nuclear geometry $\omega^*\in \Omega $ with 
    \begin{align}
        \mathcal{Q}(\omega^*,t^*)=0,
    \end{align}
    such that $(\omega^*,t^*)$ is non-degenerate in the sense of Definition  \ref{non-degen_defi}, then there exists an open set $V\subset \Omega$ with $\omega^* \in V$ and an open set $U \subset \mathbb{R}^{|\mathscr{I}_\ell|}$ with $t^*\in U$, such that there is a unique, real analytic function $t:V\longrightarrow U$ with 
    \begin{align}
        \big\lbrace (\omega,t(\omega)) \,|\,  \omega\in V\big\rbrace=\big\lbrace  (\omega,t)\in V\times U \,|\,  \mathcal{Q}(\omega,t)=0\big\rbrace.
    \end{align}
\end{theorem}

\begin{proof}
    By the regularity assumptions on the molecular orbitals and Lemma \ref{creation_annihili_analytic} the  coordinate functions $\mathcal{Q}_\upsilon(\omega,t)$ are real analytic for any $\upsilon\in \mathscr{I}_\ell$. Hence, in order to use the implicit function theorem for analytic functions, we only need to prove that the partial Jacobian 
    \begin{align}\label{jacobian_Q}
        \left[ \frac{\partial \mathcal{Q}_\upsilon (\omega^*,t^*)}{\partial t_\nu}\right]_{\upsilon,\nu=1}^{|\mathscr{I}_\ell|}
    \end{align}
    is a regular matrix. In order to prove this condition let us first calculate an arbitrary relevant partial derivative of a coordinate function, i.e.
    \begin{align*}
        \frac{\partial \mathcal{Q}_\upsilon }{\partial t_\nu}&= \frac{\partial \langle \Psi_\upsilon(\omega),\mathrm{e}^{-T(\omega,t)}\mathcal{H}_\mathfrak{M}(\omega)\mathrm{e}^{T(\omega,t)}\Psi_0(\omega)\rangle_{\mathrm{L}^2} }{\partial t_\nu}\\
        &=\big\langle \Psi_\upsilon(\omega),\big(-\mathscr{X}_\nu(\omega)\mathrm{e}^{-T(\omega,t)}\mathcal{H}_\mathfrak{M}(\omega)\mathrm{e}^{T(\omega,t)}+\mathrm{e}^{-T(\omega,t)}\mathcal{H}_\mathfrak{M}(\omega)\mathscr{X}_\nu(\omega)\mathrm{e}^{T(\omega,t)}\big)\Psi_0(\omega)\big\rangle_{\mathrm{L}^2}\\
        &=-\big\langle \Psi_\upsilon(\omega),\mathscr{X}_\nu(\omega)\mathrm{e}^{-T(\omega,t)}\mathcal{H}_\mathfrak{M}(\omega)\mathrm{e}^{T(\omega,t)}\Psi_0(\omega)\big\rangle_{\mathrm{L}^2}+\big\langle \Psi_\upsilon(\omega),\mathrm{e}^{-T(\omega,t)}\mathcal{H}_\mathfrak{M}(\omega)\mathrm{e}^{T(\omega,t)}\Psi_\nu(\omega)\big\rangle_{\mathrm{L}^2},        \end{align*}
        where we used that all excitation operators and cluster operators commute (see Proposition \ref{proposition_excitation_comm}). In the following for simplicity we will use the notation $T^*:=T(\omega^*,t^*)$. The first summand at the solution to the discretized and truncated coupled cluster equation reads
    \begin{align*}
        \langle \Psi_\upsilon(\omega^*),\mathscr{X}_\nu(\omega^*)\mathrm{e}^{-T^*}\mathcal{H}_\mathfrak{M}(\omega^*)\mathrm{e}^{T^*}\Psi_0(\omega^*)\rangle_{\mathrm{L}^2}&=\langle \mathscr{X}^{\dagger}_\nu(\omega^*)\Psi_\upsilon(\omega^*),\mathrm{e}^{-T^*}\mathcal{H}_\mathfrak{M}(\omega^*)\mathrm{e}^{T^*}\Psi_0(\omega^*)\rangle_{\mathrm{L}^2}.
        \intertext{Notice that $\mathscr{X}^{\dagger}_\nu(\omega^*)\Psi_\upsilon(\omega^*)=\mathscr{X}^\dagger_\nu(\omega^*)\mathscr{X}_\upsilon(\omega^*)\Psi_0(\omega^*)$. There are three cases to consider. Either the virtual and occupied indices of $\nu$ are contained within the occupied and virtual indices of $\upsilon$, in this case  $\mathscr{X}^\dagger_\nu(\omega^*)\mathscr{X}_\upsilon(\omega^*)=\mathscr{X}_\eta(\omega^*)$ for some excitation $\eta \in \mathscr{I}_\ell$, or $\nu=\upsilon$, in which case $\mathscr{X}^\dagger_\nu(\omega^*)\mathscr{X}_\nu(\omega^*)=\mathrm{I}$. Otherwise $\mathscr{X}^\dagger_\nu(\omega^*)\mathscr{X}_\upsilon(\omega^*)=0$. So we have}
        -\langle \Psi_\upsilon(\omega^*),\mathscr{X}_\nu(\omega^*)\mathrm{e}^{-T^*}\mathcal{H}_\mathfrak{M}(\omega^*)\mathrm{e}^{T^*}\Psi_0(\omega^*)\rangle_{\mathrm{L}^2}&=  \begin{cases}
        \,0, \\
-\langle \Psi_\eta(\omega^*),\mathrm{e}^{-T^*}\mathcal{H}_\mathfrak{M}(\omega^*)\mathrm{e}^{T^*}\Psi_0(\omega^*)\rangle_{\mathrm{L}^2}, \\
-\langle \Psi_0(\omega^*),\mathrm{e}^{-T^*}\mathcal{H}_\mathfrak{M}(\omega^*)\mathrm{e}^{T^*}\Psi_0(\omega^*)\rangle_{\mathrm{L}^2}.
\end{cases}
    \end{align*}
    Since $(\omega^*,t^*)$ is a solution, the second case also vanishes and the third case returns the (possibly truncated) coupled cluster energy $E_{\ell}$, i.e.
    \begin{align}
        \langle \Psi_\upsilon(\omega^*),-\mathscr{X}_\nu(\omega^*)\mathrm{e}^{-T^*}\mathcal{H}_\mathfrak{M}(\omega^*)\mathrm{e}^{T^*}\Psi_0(\omega^*)\rangle_{\mathrm{L}^2}=-E_{\ell}\delta_{\upsilon \nu}.
    \end{align}
    Inserting this into the derivative expression, we obtain
    \begin{align*}
        \frac{\partial \mathcal{Q}_\upsilon (\omega^*,t^*)}{\partial t_\nu}=\big\langle \Psi_\upsilon(\omega^*),\mathrm{e}^{-T^*}\mathcal{H}_\mathfrak{M}(\omega^*)\mathrm{e}^{T^*}\Psi_\nu(\omega^*)\big\rangle_{\mathrm{L}^2}-E_{\ell}\delta_{\upsilon \nu}.
    \end{align*}
So for the partial Jacobian at the solution $(\omega^*,t^*)$ we have  
\begin{align}
    \left[ \frac{\partial \mathcal{Q}_\upsilon (\omega^*,t^*)}{\partial t_\nu}\right]_{\upsilon,\nu=1}^{|\mathscr{I}_\ell|}=\mathbf{J}-E_{\ell}\cdot\mathrm{I},
\end{align}
    where $\mathbf{J}$ is the matrix with components $\mathbf{J}_{\upsilon \nu}=\big\langle \Psi_\upsilon(\omega^*),\mathrm{e}^{-T^*}\mathcal{H}_\mathfrak{M}(\omega^*)\mathrm{e}^{T^*}\Psi_\nu(\omega^*)\big\rangle_{\mathrm{L}^2}$. Therefore, the partial Jacobian of the coupled cluster function is invertible if \begin{align}\label{main_cc_det}
        \det(\mathbf{J}-E_{\ell}\cdot\mathrm{I})\ne 0.
    \end{align} 
     Since $\sigma(\mathbf{J})=\sigma(P_{\mathfrak{V}_\ell}\mathrm{e}^{-T^*}\mathcal{H}_{\mathfrak{M}}(\omega^*)\mathrm{e}^{T^*}P_{\mathfrak{V}_\ell})$, where $P_{\mathfrak{V}_\ell}$ is the orthogonal projector on the excited determinants  $\mathfrak{V}_\ell=\mathrm{span}_{\mathbb{C}}\,\mathcal{S}_\ell(\mathfrak{M}(\omega^*))\setminus  \lbrace \Psi_0(\omega^*) \rbrace$, the inequality \eqref{main_cc_det} holds by the non-degeneracy assumption.
\end{proof}

\end{document}
